\section{Problem Statement and Hypotheses}
\label{sec:problem}

\subsection{Call Graph, Faults and Outcome}

Let $G=(V,E)$ be the directed call graph of an application. $V$ is the set of services, and an
edge $(u,v)\in E$ means that $u$ calls $v$ under the traced workload. Writing
$u \rightsquigarrow v$ for a directed path from $u$ to $v$, the ancestors of a service are the
services from which it can be reached and its descendants are the services it can reach:
\begin{align}
\mathrm{anc}(v) &= \{u \in V : u \rightsquigarrow v,\ u \neq v\},\\
\mathrm{desc}(v) &= \{u \in V : v \rightsquigarrow u,\ u \neq v\}.
\end{align}

A fault $f \in \{\text{kill}, \text{latency}\}$ applied to service $v$ either terminates its
process or delays its network traffic (Section~\ref{sec:protocol}). In repetition $r$ the
degraded set $D_r(v,f)$ holds the services other than $v$ whose own latency or error rate
crosses the degradation rule of Section~\ref{sec:outcomes} during the fault window. The
degraded set splits into three disjoint parts by graph relation to the faulted service:
\begin{align}
D_r^{A} &= D_r \cap \mathrm{anc}(v), \qquad D_r^{D} = D_r \cap \mathrm{desc}(v),\\
D_r^{U} &= D_r \setminus \bigl(\mathrm{anc}(v) \cup \mathrm{desc}(v)\bigr).
\end{align}
The primary outcome of a service is the median over its repetitions of the number of degraded
ancestors,
\begin{equation}
Y(v,f) = \operatorname{median}_r \, \lvert D_r^{A}(v,f) \rvert,
\end{equation}
called the number of degraded callers below. The counts of degraded descendants and of
degraded unrelated services are recorded separately and never added to $Y$, because a combined
count would hide the direction in which a failure travels.

\subsection{Predictors}

The primary predictor is the caller count, the number of transitive callers
$X(v) = \lvert\mathrm{anc}(v)\rvert$. The comparator of interest is the hybrid criticality
metric of the original proposal,
\begin{align}
\HC(v) &= w_{\mathrm{in}} C_{\mathrm{in}}(v) + w_{\mathrm{out}} C_{\mathrm{out}}(v)\,\varphi(v)
          + w_{\mathrm{btw}} C_{\mathrm{btw}}(v),\\
\varphi(v) &= d^{+}(v) \big/ \max_{u} d^{+}(u),
\end{align}
where $C_{\mathrm{in}}$ and $C_{\mathrm{out}}$ are in-degree and out-degree centrality,
$C_{\mathrm{btw}}$ is betweenness with directed normalisation, $d^{+}$ is out-degree, and
the weights are $w_{\mathrm{in}} = \val{hybrid_w_in}$, $w_{\mathrm{out}} = \val{hybrid_w_out}$
and $w_{\mathrm{btw}} = \val{hybrid_w_btw}$ (Section~\ref{sec:predictors}). The factor
$\varphi$ grows linearly with out-degree, so the metric rewards services that call many others.
Ten further graph measures are kept as comparators. Each test pairs one predictor with one
outcome across the services of one application, with the gateway left out
(Section~\ref{sec:analysis-set}).

\subsection{Hypotheses}

The hypotheses are those of the locked preregistration, where the second carries the label H1b
because the name H2 there belongs to an untested structural observation. H1, the primary
hypothesis, states that the caller count predicts $Y$ while out-degree, descendant count and
the hybrid metric do not. It is tested as a positive Spearman rank correlation:
\begin{equation}
\text{H1:}\quad \rho(X, Y) > 0.
\end{equation}
H2 is directional. It states that the fan-out premise points the wrong way, so that the hybrid
metric is not positively associated with the outcome:
\begin{equation}
\text{H2:}\quad \rho(\HC, Y) \leq 0.
\end{equation}
The null hypothesis of each test is that predictor and outcome have no monotone association.
Every p-value is two-sided, comes from a permutation test and is corrected with the
Benjamini-Hochberg procedure within its family (Section~\ref{sec:stats}). Kill faults form the
confirmatory family, and latency faults repeat every test as a separate replication family.
